% Generated draft: requires mathematical and source review.
\documentclass[11pt]{article}
\usepackage[margin=25mm]{geometry}
\usepackage{fontspec}
\setmainfont{Latin Modern Roman}
\usepackage{amsmath,amssymb,mathtools,unicode-math}
\setmathfont{Latin Modern Math}
\usepackage{enumitem,needspace,xurl,hyperref}
\setlength{\parindent}{0pt}
\setlength{\parskip}{5pt}
\setlength{\emergencystretch}{3em}
\newcommand{\N}{\mathbb{N}}
\newcommand{\Z}{\mathbb{Z}}
\newcommand{\Q}{\mathbb{Q}}
\newcommand{\R}{\mathbb{R}}
\newcommand{\C}{\mathbb{C}}
\newcommand{\F}{\mathbb{F}}
\newcommand{\A}{\mathbb{A}}
\newcommand{\PP}{\mathbb{P}}
\newcommand{\E}{\mathbb{E}}
\newcommand{\eps}{\varepsilon}
\newcommand{\dd}{\,\mathrm{d}}
\begin{document}
\section*{UnsolvedMath ID: 10600046}
\textbf{Source problem code:} catalogue entry
\subsection{Definitions and mathematical statement}
In an oriented long virtual knot diagram, traverse from the incoming end. Each classical crossing is early-over or early-under according as its first encountered branch passes over or under. The source describes an invariant method assigning two related but distinct quandle-type operations to these two crossing classes. It reports detection of non-invertibility and of noncommutativity under long-knot concatenation, but very little information for classical knots and failure of this construction at the Kauffman-bracket-polynomial level.

The item asks for two extensions.
\begin{enumerate}
\item Carry the two-crossing-class/two-operation method to other algebraic invariant constructions, including what the source calls nonlinear quandles and quandles over rings with non-unique decomposition of multiples. Find compatibility laws that make the structure invariant while retaining enough distinction to detect reversal or noncommutativity. The quoted ring and nonlinearity terminology is not defined and is not replaced by a specific factorization hypothesis.
\item On the usual Khovanov chain space \(C\), construct two different differentials \(d_\circ,d_\star\) corresponding to the two operations, with
\[
 d_\circ^2=d_\star^2=0,\qquad
 d_\circ d_\star=d_\star d_\circ.
\]
Specify an invariant construction from this data and investigate its ability to detect non-invertibility. A suggested source of the two differentials is ordinary and odd Khovanov theory, after an identification \(\phi:C_{\mathrm{odd}}\to C_{\mathrm{ordinary}}\), possibly noncanonical: one could seek \(d_\circ=d_{\mathrm{ordinary}}\) and \(d_\star=\phi d_{\mathrm{odd}}\phi^{-1}\), provided the required common-space and commuting properties can actually be established.
\end{enumerate}
The source explicitly requests commuting differentials; this has not been replaced by anticommutation. In particular, square-zero of their sum is not automatic with this convention outside characteristic \(2\). Coefficients, gradings, the identification's choice dependence, Reidemeister comparison data and the resulting invariant are unspecified. Detection is an intended application, not a claim that every noninvertible knot will be detected. The reversal and concatenation conventions and the original two-operation construction are also absent.
\subsection{Short English statement}
Extend the early-over/early-under two-operation method to other algebraic invariants, and construct a Khovanov-level invariant with two distinct commuting differentials, aiming to detect knot reversal.
\subsection{Sources}
\begin{itemize}
\item \emph{Unsolved Problems in Virtual Knot Theory and Combinatorial Knot Theory}, revised survey, problem-list item 46. \url{https://arxiv.org/abs/1409.2823}. Recorded access 2026-07-29; extraction source-tex. The initial reference for the two-operation method is missing.
\item The supplied summary lists \emph{On the non-invertibility of virtual knots} without full citation, \emph{Vassiliev Invariants from Parity Mappings} (2012, arXiv:1203.xxxx), \emph{Free Knots, Groups, and Finite-Type Invariants} (2010, arXiv:1004.xxxx), and \emph{Virtual Covers of Links} (2014, arXiv:1405.xxxx), as well as unspecified chirality studies. The xxxx identifiers are incomplete placeholders.
\item Local UnsolvedMath registry, \texttt{data/unsolvedmath/problems.json}, permanent ID \(10600046\), exact source code \texttt{catalogue entry}; supplied partially\_solved/PARTIAL-PROGRESS status.
\end{itemize}
\end{document}
